\section{结论}

本文基于统一办公建筑原型，建立了适用于五个代表性城市工况的空调方案比较框架，并在同一边界条件下完成了理想负荷、VRF+DOAS和FCU+DOAS三层结果的系统分析。通过先建立理想负荷基线，再引入两类实际空调系统，本文将气候驱动建筑负荷与系统侧运行表现进行了相对清晰的分离，为课程设计中的技术比较提供了可复现基础。

结果表明，深圳是制冷需求最强的城市，其峰值冷负荷达到\SI{327.092}{\kilo\watt}，年冷负荷也最高；沈阳则具有最高的年热负荷。重庆作为新增的夏热冬冷城市，与成都共同补充了中西部湿热地区的比较样本。设备选型方面，设计冷量需求总体上随制冷强度增加而增加，而在当前仅统计空调电耗的比较口径下，FCU+DOAS在五个城市中均表现出更低的年空调电耗强度，但这一结论尚不等同于完整系统综合能耗更优，因为锅炉燃料消耗尚未纳入统计。

因此，本文更适合作为不同气候分区下空调方案适宜性筛选和方法比较的阶段性成果，而非最终工程优选结论。后续研究应结合GB~50189-2015进一步引入围护结构分区化设计，同时补齐FCU+DOAS的燃料消耗统计，并将评价指标扩展到综合能耗、碳排放和控制策略层面，以形成更完整的工程判断依据。
